\documentclass[11pt]{amsart}
\usepackage[T1]{fontenc}
\usepackage[utf8]{inputenc}
\usepackage{amssymb}
\usepackage[margin=1.05in]{geometry}
\usepackage[hidelinks]{hyperref}
\newcommand{\stamptime}{\the\numexpr\time/60\relax:\ifnum\numexpr\time-60*(\time/60)\relax<10 0\fi\the\numexpr\time-60*(\time/60)\relax}
\AddToHook{shipout/foreground}{\put(0.5\paperwidth,-\paperheight+0.55cm){\makebox[0pt]{\tiny\ttfamily unit/doc/\jobname.tex \ --- \ compiled \today, \stamptime}}}
\newtheorem{theorem}{Theorem}
\newtheorem{proposition}[theorem]{Proposition}
\newtheorem{lemma}[theorem]{Lemma}
\newtheorem{corollary}[theorem]{Corollary}
\theoremstyle{definition}
\newtheorem{definition}[theorem]{Definition}
\theoremstyle{remark}
\newtheorem{remark}[theorem]{Remark}
\newcommand{\Om}{\Omega}
\newcommand{\Map}{\mathrm{Map}_*}
\newcommand{\Z}{\mathbb{Z}}
\newcommand{\code}[1]{\texttt{#1}}
\newcommand{\cT}{\mathsf{isTrunc}}
\newcommand{\cC}{\mathsf{isConn}}
\newcommand{\refl}{\mathsf{refl}}
\newcommand{\lp}{\mathsf{loop}}
\newcommand{\Kf}{\mathsf{Kf}}
\newcommand{\Hom}{\mathrm{Hom}}
\newcommand{\hc}[2]{[#1,#2]_*}
\newcommand{\Unl}{\mathsf{Unloop}}
\newcommand{\PMO}{\mathsf{PMO}}
\newcommand{\Fib}{\mathrm{Fib}}
\newcommand{\kb}{\bar\kappa}
\newcommand{\sq}{\mathsf{sq}}
\newcommand{\cE}{\mathcal E}
\title{The cross term of the reduced cup square, and a delooping of $\Fib\bar\kappa$}
\date{28 October 2026}
\begin{document}
\maketitle

\noindent\textbf{Purpose and status.} First task of \S2 of
\code{doc/pi3-s2-postnikov-handoff-9.txt} (towards $\pi_4(S^3)=\Z/2$):
an informal proof that the cross term of $\kb=\kappa;K(\rho,4)$
vanishes, and that $R=\Fib\kb$ consequently deloops, with a homogeneous
structure and a torsor as corollaries. Written before any code; the
numbered statements are meant to be declarations (\S\ref{s:corr}).
Foundations as in the development (univalence, images; function
extensionality derived). No new axiom; \emph{the $S^3$ axiom is not used}
(it is needed only for step~3 of the handoff); no theory of quadratic
maps, no pushout, no smash product. Refereed twice (28 Oct; second review of the revision: no further
obstruction): main route sound; repairs incorporated
(Lemma~\ref{l:pi1} for $\PMO_1$ with arbitrary $A$, in place of
\code{PICircle}; $\tau$ as an explicit pointed equivalence through
\code{PtMapLoops}, Lemma~\ref{l:tau}; the converse in \S7 needs a
normalisation of the cell and is withdrawn as open; Lemma~\ref{l:E2}
with the locks of $e_1$ and of the circle generator made explicit;
citations corrected).

\smallskip\noindent\textbf{The idea.} Write $K_m=K(\Z,m)$, $K'_m=K(\Z/2,m)$,
$X=K_3$, $W=K'_5$, and let $f=e_2;\kb;e_4^{-1}:\Om X\to_*\Om W$ be $\kb$
seen on loop spaces. By DW1's Lemma~9 (\code{DeloopMaps}), $f$ is the
looping of a map $\hat F:X\to_*W$ as soon as the cell at the base point
is inhabited, and at $c=2$, $t=1$ the cells are propositions. So
\emph{no pointing in $p$ is needed}: it suffices to have, for every
$p:\Om X$, a path $f=D_p$ in $\PMO$, where $D_p(q)=(fp)^{-1}\cdot f(p\cdot q)$.
These paths come from the \emph{cross term}, the bipointed map
$\tilde\delta(x)(y)$ ``$=\kb(x+y)-\kb(x)-\kb(y)$'', which we show is
$0$ by Theorem~A of \code{doc/h4.pdf}: its diagonal is
``$\kb(2x)-2\kb(x)$'', and both terms vanish \emph{by construction} (not as
classes) because $2=0$ in $\Z/2$. Then $BR=\Fib\hat F$ is a pointed
2-connected 4-type with $\Om BR\simeq_*R$.

\smallskip\noindent\textbf{Where $2=0$ enters, twice.} (i) Doubling on
$\Om K'_{m+1}$ is $e_m;K(2,m);e_m^{-1}$ (Lemma~\ref{l:dbl}), and
$K(2,m)=K(0,m)=0$ over $\Z/2$. (ii) Precomposing $c_2$ in the first
variable with $K(2,2)$ multiplies its invariant $\cE$ by $2$
(Lemma~\ref{l:E2}); after $K(\rho,4)_*$ this is again doubling in
$\Om^2K'_4$, which is $0$ (Theorem~\ref{t:kill}). Over $\Z$ neither
vanishes (the general formula, Proposition~\ref{p:formula}, gives the
cross term $c_2;K(2g,4)_*$).

\section{Conventions}

As in \code{doc/h4.pdf}: \code{isTrunc n} is ``$(n-2)$-truncated'',
\code{isConn n} is ``$(n-2)$-connected''; $K(A,m)=\code{KK[A].K m}$ with
$e_m:\Om K(A,m{+}1)\simeq_*K(A,m)$ (\code{KK.e}, $e_0$ the band); $F;G$ is
\code{ptcomp F G}; $p\cdot q$ is \code{trans p q} (so
$p\cdot\refl\equiv p$, $\refl\cdot\refl\equiv\refl$); for pointed $X,Y$,
$\Map(X,Y)$ is pointed by $0=(\lambda\_.y_0,\refl)$; $M_Y=\Map(K_2,Y)$;
$\lp$, $D$, $\Kf$, $E$, $\Unl$, $\mathsf{Rel}$, $\mathsf{Levels}$ are those of
\code{DeloopMaps.Maps}; $\PMO_m=\Map(\Om K(A,m{+}1),\Om K(A,m{+}1))$;
$\mathrm{conj}_m F=e_m;F;e_m^{-1}$ (\code{KMap.conj}), so
$\lp K(h,m{+}1)=\mathrm{conj}_m K(h,m)$ (\code{loopK}). From
\code{doc/h4.pdf}: $c_2:K_2\to_*M_{K_4}$, $\kappa=\Delta c_2$ (by
definition), $\cE$ and $\Theta_B=\cE;j_B$ (\code{H4Bilinear.EDef},
\code{BilAt.Theta}), Theorem~A, Theorem~B, Lemma~j, the pointing
irrelevance lemma \code{pi-set}.

\section{The group $\Z/2$ and the data}

\begin{definition}\label{d:z2}
$\Z/2$ is the abelian group (\code{Torsor.Ab}) with carrier \code{Bool},
$o=\mathsf{false}$, $a+b=\mathsf{xor}(a,b)$, $-a=a$; it is a set
(decidable equality, \code{Hedberg}); the laws hold by cases. For an
abelian $A$, $2_A:\Hom(A,A)$ is $a\mapsto a+a$ (associativity and
commutativity), $0_{A,B}:\Hom(A,B)$ is constant $o$. The reduction
$\rho:\Hom(\Z,\Z/2)$ is the parity, with $\rho(\mathsf{suc}\,x)=\rho(\mathsf{pred}\,x)=\mathsf{not}(\rho\,x)$,
so $\rho(x+y)=\rho x+\rho y$ by induction on $y$, as in \code{Integers}.
\end{definition}

\begin{lemma}\label{l:two}
$2_{\Z/2}=0_{\Z/2,\Z/2}$ in $\Hom(\Z/2,\Z/2)$.
\end{lemma}
\begin{proof}
$\mathsf{xor}(a,a)=\mathsf{false}$ by cases; the homomorphism parts agree
because the carrier is a set.
\end{proof}

\noindent Let $\kb=\kappa;K(\rho,4):K_2\to_*K'_4$ and $R=\Fib\kb$.
By \code{PostnikovRealization.Triple} at $(\Z,\Z/2,\kb)$, $R$ is
\code{isConn 3}, \code{isTrunc 5}, and $\pi_3R\cong\Z/2$ as groups.

\section{Doubling on Eilenberg--MacLane loop spaces}

For a pointed type $(Y,y_0)$ let $\sq:\Om Y\to_*\Om Y$, $\sq(p)=p\cdot p$,
pointed by $\refl$ ($\refl\cdot\refl\equiv\refl$).

\begin{lemma}[pointing irrelevance, 1-types]\label{l:pi1}
Let $(P,p_0)$ be \code{isConn 2} and $(Q,q_0)$ \code{isTrunc 3} with
$\Om(Q,q_0)$ commutative. Then $(g,\zeta)=(g,\zeta')$ for all $g:P\to Q$
and $\zeta,\zeta':g\,p_0=q_0$.
\end{lemma}
\begin{proof}
By \code{unique-loop1} ($c=0$; $3=1+(1+1)$) it suffices that
$\lp(g,\zeta)=\lp(g,\zeta')$ in $\Map(\Om P,\Om Q)$. By \code{loop-ap} the
underlying maps are $q\mapsto\zeta^{-1}\cdot\mathsf{ap}_gq\cdot\zeta$ and the
same with $\zeta'$; they differ by conjugation by the loop
$\zeta^{-1}\cdot\zeta'$ of $Q$ (\code{LoopConj.conj-trans}), which is trivial
since $\Om Q$ is commutative. The pointings lie in identity types of the
set $\Om Q$, which are propositions.
\end{proof}

\begin{lemma}[pointing irrelevance on $\PMO_m$]\label{l:pmo-pi}
For every $m$, $\PMO_m$ is a set, and two elements of $\PMO_m$ with
pointwise equal underlying maps are equal.
\end{lemma}
\begin{proof}
Set: $m=0$, the target is the set $\Om K(A,1)$; $m\ge1$, \code{cor5-map}
(source \code{isConn (m+1)}, target \code{isTrunc (m+2)}). Pointings: two
pointings of one map differ by a loop of $\Om K(A,m{+}1)$ at $\refl$,
i.e.\ an element of $\Om^2K(A,m{+}1)$. $m=0$: this type is a proposition.
$m=1$: Lemma~\ref{l:pi1} with $P=Q=\Om K(A,2)$ (2-loops commute,
\code{EH.eh}); \code{PICircle} is not used, it is the circle property of
$K(\Z,1)$ only. $m\ge2$: $\Om^2K(A,m{+}1)$ is connected
($\simeq K(A,m{-}1)$) and \code{pi-set} applies.
\end{proof}

\begin{lemma}[$K$ of zero]\label{l:K0}
$K(0_{A,B},m)=0$ for all $m$.
\end{lemma}
\begin{proof}
$m=0$: both underlying maps are constant $o$; pointings in a set.
$m+1$: by \code{KFun.Unloop.unloop} it suffices that the loopings agree:
$\lp K(0,m{+}1)=\mathrm{conj}_m K(0,m)=\mathrm{conj}_m 0=0=\lp 0$
(\code{loopK}, induction, composition with a constant map, \code{loop-const}).
\end{proof}

\begin{lemma}[doubling]\label{l:dbl}
For every abelian $A$ and $m$, $\sq=\mathrm{conj}_m K(2_A,m)$ in $\PMO_m$,
i.e.\ $\sq=\lp K(2_A,m{+}1)$.
\end{lemma}
\begin{proof}
By Lemma~\ref{l:pmo-pi} it suffices to compare underlying maps, by
induction on $m$.
$m=0$: $\mathrm{conj}_0K(2,0)=\mathsf{unband}\circ2\circ\mathsf{band}$
(\code{c0-conj}), and $\mathsf{band}(p\cdot p)=\mathsf{band}\,p+\mathsf{band}\,p$
(\code{band-hom}) and $\mathsf{band}(\mathsf{unband}\,a)=a$ (\code{band-unband}); conclude with \code{band-inj}.
$m+1$: let $G=e_{m+1}^{-1};\sq;e_{m+1}:K(A,m{+}1)\to_*K(A,m{+}1)$. It
suffices that $G=K(2,m{+}1)$ (then $\sq=e_{m+1};G;e_{m+1}^{-1}$ by
\code{pt-inv-l/r}), hence, by \code{KFun.Unloop.unloop}, that
$\lp G=\lp K(2,m{+}1)=\mathrm{conj}_mK(2,m)$, an equality in $\PMO_m$,
so pointwise (Lemma~\ref{l:pmo-pi}). By \code{loop-comp},
$\lp G=\lp e_{m+1}^{-1};\lp\sq;\lp e_{m+1}$. (a) $\lp\sq$ is $\sq$ on
$\Om^2K(A,m{+}2)$ pointwise: $(\lp\sq)(\alpha)=\refl^{-1}\cdot\mathsf{ap}_\sq\alpha$
and $\mathsf{ap}_{q\mapsto q\cdot q}\alpha=\alpha\cdot\alpha$ for
$\alpha:\refl=\refl$ (a new lemma \code{ap-diag}: $\mathsf{ap}$ of $q\mapsto q\cdot q$ along $a:p=p'$ is
$\mathsf{ap}_{-\cdot p}a\cdot\mathsf{ap}_{p'\cdot-}a$, by $J$; then \code{EH.whisk-r},
\code{EH.whisk-l}; \code{Groupoid.ap-trans} is a different statement).
(b) With $\gamma=(\lp e_{m+1}^{-1})\beta$:
$(\lp e_{m+1})(\gamma\cdot\gamma)=(\lp e_{m+1})\gamma\cdot(\lp e_{m+1})\gamma$
(\code{loop-mult}) and $(\lp e_{m+1})\gamma=\beta$ (\code{loop-comp},
\code{pt-inv-r}, \code{loop-id}). So $(\lp G)\beta=\beta\cdot\beta=\sq\,\beta$,
which is $(\mathrm{conj}_mK(2,m))\beta$ by induction.
\end{proof}

\begin{corollary}\label{c:sq0}
Over $\Z/2$: $\sq=0$ in $\PMO_m$ for every $m$; so there are homotopies
$s_m:\prod_{p:\Om K'_{m+1}}p\cdot p=\refl$. Also, for
$w:\Om^2K'_4$, $w\cdot w=\refl$ (a homotopy $s'$).
\end{corollary}
\begin{proof}
Lemma~\ref{l:dbl}, Lemma~\ref{l:two} (\code{apH} of $K(-,m)$), Lemma~\ref{l:K0},
and $\mathrm{conj}_m0=0$. For $s'$: $\lp e_3:\Om^2K'_4\to_*\Om K'_3$ is an
equivalence with $(\lp e_3)(w\cdot w)=((\lp e_3)w)^{2}$ (\code{loop-mult});
apply $s_2$; $(\lp e_3)\refl=\refl$ is the pointing of $\lp e_3$, and
injectivity of $\lp e_3$ comes from its being an equivalence (\code{LoopIso.loopPE}).
\end{proof}

\section{Precomposition with $K(2,2)$}

\begin{lemma}\label{l:E2}
For pointed $Y$ and $c:K_2\to_*M_Y$, the underlying map of
$\cE(K(2_\Z,2);c)$ is pointwise $x\mapsto\cE(c)(x)\cdot\cE(c)(x)$.
\end{lemma}
\begin{proof}
$\cE(c)=\mathsf{ev}(\mathsf{evPt}(e_1^{-1};\lp c;\mathsf{ev}_Y))$, with
$e_1=\code{CupSquare.e1}$, which is $\code{KK[Zab].e 1}$ \emph{by definition}
but locked: a new lemma \code{e1-def} $:e_1=\code{KZ.e (suc zero)}$ by $\refl$
is added to \code{CupSquare} just before \code{opaque e1} (rule~3; no
statement changes). By \code{loop-comp} and \code{KMap.loopK 1} (raw
\code{KK.e}, not the locked \code{Lev.eA}),
$\lp(K(2,2);c)=e_1;K(2,1);e_1^{-1};\lp c$; with
$L=e_1^{-1};\lp c;\mathsf{ev}_Y$ the argument of $\mathsf{evPt}$ is
$K(2,1);L$ (\code{e1-def}, \code{ptcomp-assoc}, \code{pt-inv-r},
\code{ptcomp-idl}). Let $g_N$ be the generator of the instance
\code{Circle[N]}, $N=\Map(K_2,\Om Y)$, used by $\mathsf{evPt}$ in $\cE$
(each instance locks its own $g$; compare \code{ENat.g-eq}). Then
$\mathsf{evPt}(K(2,1);L)\equiv(\lp(K(2,1);L))(g_N)=(\lp L)((\lp K(2,1))\,g_N)$
(\code{evPt-loop}, \code{loop-comp}), and
$(\lp K(2,1))\,g_N=g_N\cdot g_N$ by \code{loopK 0} and Lemma~\ref{l:dbl}
at $m=0$, $A=\Z$, applied at $g_N$ (not through $\mathsf{evPt}$ of a different
instance; formally \code{k21-sq}, with $K(2,1)$ written once at the types
of the circle, \code{k21}, where $K(\Z,1)$ of the circle is identified
with \code{KK.K 1}). By
\code{loop-mult} it is $\mathsf{evPt}\,L\cdot\mathsf{evPt}\,L$, and $\mathsf{ev}$ takes
a product of loops of $N$ to the pointwise product (\code{ev-ap},
\code{ap-trans}).
\end{proof}

\begin{lemma}\label{l:O2mult}
For $G:Y\to_*Y'$, $(\lp^2G)(\alpha\cdot\beta)=(\lp^2G)\alpha\cdot(\lp^2G)\beta$
(\code{loop-mult} for $\lp G$).
\end{lemma}

\begin{theorem}[the reduced cup square kills doubling]\label{t:kill}
$c''=K(2,2);c_2;K(\rho,4)_*$ is $0$ in $\Map(K_2,M_{K'_4})$. Hence there
is a homotopy $N:\prod_{x:K_2}\kb(K(2,2)\,x)=*$.
\end{theorem}
\begin{proof}
$\Theta_{\Z/2}$ is injective (\code{BilAt.theta-inj}), and pointed maps
$K_2\to_*K'_2$ with equal underlying maps are equal (\code{pi-set}). Pointwise,
$\cE(c'')=\cE(K(2,2);c_2);\lp^2K(\rho,4)$ (\code{ENat.e-nat}, the
$\mathsf{ptcomp}$ of the two outer maps is the one of the definition),
which is $x\mapsto(\cE(c_2)x\cdot\cE(c_2)x)$ mapped by $\lp^2K(\rho,4)$
(Lemma~\ref{l:E2}, $\cE(c_2)=\iota$ by \code{cup-norm}), i.e.\
$x\mapsto w_x\cdot w_x$ with $w_x=(\lp^2K(\rho,4))(\iota x)$
(Lemma~\ref{l:O2mult}); this is $\refl$ by $s'$
(Corollary~\ref{c:sq0}). Also $\cE(0)$ is pointwise $\refl$: \code{loop-const} and
\code{comp-const-l/r} give $\mathsf{phi}\,0=0$, then \code{evPt-loop} with
\code{loop-const}, and $\mathsf{ev}\,\refl\equiv r$. $\Theta(c'')$ and $\Theta(0)$ are then
pointwise the point through the pointing of $j_B$, and \code{pi-set}
(\code{BilAt.sK2}, \code{cK2}) makes them equal, as in \code{key}. The
same $K(\rho,4)$, the locked \code{Lev[Zab,Z2ab].K rho four} (as in
\code{BilAt.c2g}), is used both in $\kb$ and in $K(\rho,4)_*$, so that the
unfolding for $N$ below is definitional (with the raw \code{KMap.K four} the
checker unfolded its construction, \code{four} being a numeral). So
$\Theta(c'')=\Theta(0)$ and $c''=0$.
For $N$: by definition $c''(x)=c_2(K(2,2)x);K(\rho,4)$, so
$c''(x)_1(K(2,2)x)\equiv K(\rho,4)(\kappa(K(2,2)x))\equiv\kb(K(2,2)x)$; evaluate
$c''=0$ at $x$ and then at $K(2,2)x$.
\end{proof}

\section{The cross term, and why it vanishes}

From now on $X=K_3$, $W=K'_5$ (pointed by $x_0$, $w_0$), with
$\cC\,4\,X$ and $\cT\,7\,W$. Let
$f=e_2;\kb;e_4^{-1}:\Om X\to_*\Om W$ (\code{KMap.conj} shape), an element
of $\PMO$ of $\code{Maps[X,x0,W,w0]}$, and $D_p=D\,x_0\,w_0\,f_1\,p:\PMO$
for $p:\Om X$, so $D_p(q)=(f p)^{-1}\cdot f(p\cdot q)$ by definition.

\begin{lemma}[translation in $\PMO$]\label{l:tau}
Let $L=\code{PtMapLoops[}\Om X,\refl,W,w_0\code{]}$, so
$\mathsf{ev}:\Om\Map(\Om X,W)\simeq_*\PMO$ with $\mathsf{ev}\,\refl\equiv0$
(\code{ev-pt}), inverse $\mathsf{unev}$ (\code{ev-unev}). Put
$\ell=\mathsf{unev}\,f$ and
\[\tau(h)=\mathsf{ev}(\mathsf{unev}\,h\cdot\ell^{-1}),\qquad
\tau_0=\mathsf{ap}_{\mathsf{ev}}(\code{trans-sym-r}\ \ell):\tau f=0 .\]
Then $(\tau,\tau_0):(\PMO,f)\simeq_*(\PMO,0)$ (a composite of equivalences:
$\mathsf{unev}$, right translation by $\ell^{-1}$ with inverse right
translation by $\ell$, $\mathsf{ev}$), and pointwise
$\tau(h)_1(q)=h_1(q)\cdot(f_1q)^{-1}$ (\code{ev-ap}: $\mathsf{ev}$ of a path
is its pointwise $\mathsf{ap}$; $\mathsf{ap}$ of a composite and of an
inverse; $\mathsf{ev}\circ\mathsf{unev}=\mathrm{id}$). The pointing
coherence is part of the construction, not discarded: $\PMO$ is a
2-type. Only the pointwise formula is used later (Lemma~\ref{l:diag}).
\end{lemma}

\begin{definition}[the cross term]\label{d:cross}
$\tilde\delta=e_2^{-1};\hat D;\tau;\mathrm{Cj}:K_2\to_*M_{K'_4}$, the composite of
\begin{itemize}
\item $\hat D=(p\mapsto D_p,\ (E f)^{-1}):(\Om X,\refl)\to_*(\PMO,f)$;
\item $\tau:(\PMO,f)\simeq_*(\PMO,0)$, Lemma~\ref{l:tau};
\item $\mathrm{Cj}:(\PMO,0)\simeq_*(M_{K'_4},0)$, $h\mapsto e_2^{-1};h;e_4$
  (inverse $e_2;-;e_4^{-1}$; \code{comp-const-l/r}).
\end{itemize}
So, by Lemma~\ref{l:tau} for the middle factor, $\tilde\delta(x)(y)=e_4\bigl((f\bar x)^{-1}\cdot f(\bar x\cdot\bar y)\cdot(f\bar y)^{-1}\bigr)$
with $\bar x=e_2^{-1}x$: ``$\kb(x+y)-\kb x-\kb y$''.
\end{definition}

\begin{lemma}\label{l:cell}
If $\tilde\delta=0$ then $f=D_p$ for every $p:\Om X$.
\end{lemma}
\begin{proof}
At $x=e_2p$, $\tilde\delta(x)=0=\mathrm{Cj}(\tau f)$ (pointings of $\tau$,
$\mathrm{Cj}$); $\tau$ is injective as an equivalence of $\PMO$ (Lemma~\ref{l:tau};
$\PMO$ is not a set, so \code{pi-set} cannot be used), and $\mathrm{Cj}$ by
\code{pre-inj}, \code{post-inj}; so $D_{\bar x}=f$, and $\bar x=e_2^{-1}(e_2p)=p$
(\code{sec-ret}).
\end{proof}

\begin{lemma}[the diagonal of the cross term]\label{l:diag}
$\Delta\tilde\delta$ is pointwise equal to the constant map; hence
$\hc{\Delta\tilde\delta}{}=[0]$ in $\hc{K_2}{K'_4}$.
\end{lemma}
\begin{proof}
By the pointwise formula of Lemma~\ref{l:tau} (not by unfolding: $\tau$ is
built through $\mathsf{ev}/\mathsf{unev}$), $\Delta\tilde\delta(x)=e_4\bigl((fp)^{-1}\cdot f(p\cdot p)\cdot(fp)^{-1}\bigr)$,
$p=\bar x$; the rest of $\tilde\delta$ ($e_2^{-1}$, $\hat D$, $\mathrm{Cj}$) unfolds by definition. By Lemma~\ref{l:dbl} ($A=\Z$, $m=2$),
$p\cdot p=e_2^{-1}(K(2,2)(e_2p))$, so
$f(p\cdot p)=e_4^{-1}(\kb(K(2,2)(e_2p)))$ ($e_2(e_2^{-1}z)=z$, the \code{sec-eq} of $e_2$),
$=e_4^{-1}(*)=\refl$ ($N$, Theorem~\ref{t:kill}; pointing of $e_4^{-1}$). Then
$(fp)^{-1}\cdot\refl\cdot(fp)^{-1}\equiv\sq((fp)^{-1})=\refl$ ($s_4$,
Corollary~\ref{c:sq0}), and $e_4(\refl)=*$. For the class: a pointed map
$(k,\zeta)$ with $H:k\sim\mathrm{const}$ equals $(\mathrm{const},H(*)^{-1}\cdot\zeta)$
(\code{ptmap-path}), and $H(*)^{-1}\cdot\zeta$ is a loop of $K'_4$, merely
$\refl$ ($K'_4$ is \code{isConn 5}); the goal is a proposition.
\end{proof}

\begin{theorem}[the cross term vanishes]\label{t:cross}
$\tilde\delta=0$; so $f=D_p$ for every $p$.
\end{theorem}
\begin{proof}
By Theorem~A for $B=\Z/2$, $c\mapsto[\Delta c]$ is injective on
$\Map(K_2,M_{K'_4})$ (\code{DiagAt.diag-inj} and \code{ClassesAt.cls-inj};
$\mathsf{cls}(\Delta'c)=[\Delta c]$ by definition). $\Delta0\equiv0$, and
$[\Delta\tilde\delta]=[0]$ by Lemma~\ref{l:diag}. Then Lemma~\ref{l:cell}.
\end{proof}

\section{The delooping}

\begin{theorem}\label{t:unloop}
$\Unl f$ is contractible.
\end{theorem}
\begin{proof}
As \code{Over.Mult.cor12}, with the cell given directly. Take
$\mathsf{Levels}[c{=}2,t{=}1]$ ($\cC\,4\,X$; $\cT\,7\,W$ and
$7=3+(3+1)$): each $(b:W)\times C\,a\,b$ is a proposition (\code{cell}) and
$\Unl f$ is a proposition (\code{cor10}). At $a=x_0$,
$(w_0,f_1,\lambda p.\,\text{Theorem~\ref{t:cross}})$ inhabits it, so every cell
is inhabited (\code{conn-elim 2}; \code{Mult} uses \code{conn-elim zero}, $c=0$), giving $\mathsf{rel}:\mathsf{Rel}$, and
$\code{lemma9.r}\ \mathsf{rel}:\Unl f$.
\end{proof}

\begin{definition}\label{d:BR}
$(\hat F,\gamma)$ is the centre: $\hat F:K_3\to_*K'_5$, $\gamma:f=\lp\hat F$.
$BR=\Fib\hat F$, pointed by $(x_0,\hat F_0)$.
\end{definition}

\begin{theorem}[$R$ deloops]\label{t:BR}
(i) $\Om BR\simeq_*R$. (ii) $BR$ is \code{isTrunc 6} and \code{isConn 4}
(a pointed 2-connected 4-type).
\end{theorem}
\begin{proof}
(i) $\Om\Fib\hat F\simeq_*\Fib(\lp\hat F)$ (\code{Fibres.omega-fib}),
$=\Fib f$ along $\gamma^{-1}$ (\code{fib-path}), $\simeq_*\Fib\kb=R$ along
$e_2$, $e_4$ (\code{fib-equiv}; $f=e_2;\kb;e_4^{-1}$).
(ii) Truncation: $X$ is \code{isTrunc 5}, hence \code{isTrunc 6}
(\code{trunc-suc}), and the identity types of $W$ are \code{isTrunc 6}. Connectedness: $BR$ is 0-connected (\code{FibConn}, with the
surjectivity on loops as in \code{surj-triv}), and $\Om BR\simeq R$ is
\code{isConn 3}; \code{conn-from-point 2} gives \code{isConn 4}.
\end{proof}

\begin{corollary}[homogeneous structure, torsor]\label{c:mu}
$R$ has a homogeneous structure $\mu:(r:R)\to R\simeq_*(R,r)$, hence a
torsor $\mathsf{homog}\,\mu:\mathsf{Tor}_R$ (\code{Deloop.Torsors[R,r0]}),
so $\mathsf{PT}$ is contractible (\code{cor8}, $c=1$: $R$ is \code{isConn 3},
\code{isTrunc 5}$=$\code{isTrunc (2+(2+1))}), and $\mathsf{Tor}_R$ is a
connected delooping of $R$ (\code{theorem2}, \code{tor-conn},
\code{tor-trunc}).
\end{corollary}
\begin{proof}
On $\Om BR$, $y\mapsto(q\mapsto y\cdot q)$ is a pointed equivalence
$(\Om BR,\refl)\simeq_*(\Om BR,y)$ ($y\cdot\refl\equiv y$; inverse
$y^{-1}\cdot-$). Transport along (i): for $\phi:R\simeq_*\Om BR$,
$\mu(r)=\phi;(\phi r\cdot-);\phi^{-1}$, pointed by \code{ret-eq}.
\end{proof}

\begin{remark}[direct homogeneity]
Corollary~\ref{c:mu} can also be read off Theorem~\ref{t:cross} without
$\hat F$: on $\Fib f$, $T_{(p,\alpha)}(q,\beta)=(p\cdot q,\ \mu_{p,q}\cdot\mathsf{ap}_2(\alpha,\beta))$,
with $\mu_{p,q}:f(p\cdot q)=fp\cdot fq$ from $f=D_p$, and $T(\refl,f_0)=(p,\alpha)$
from the pointing part of $f=D_p$. Only the pointing in $q$ is used;
the delooping above uses even less (no pointing in $p$). We use
Theorem~\ref{t:BR}, which also gives $BR$ as a fibre $K(\Z,3)\to K(\Z/2,5)$,
a presentation suited to step~3.
\end{remark}

\section{The general formula (not used)}

\begin{proposition}[cross-term formula]\label{p:formula}
For $g:\Hom(\Z,B)$ and $F=\kappa;K(g,4)$, the cross term $\tilde\delta_F$
(Definition~\ref{d:cross} with $B$ for $\Z/2$) is $c_2;K(2g,4)_*$. Hence
$e_2;F;e_4^{-1}$ deloops if $2g=0$. (The converse, which would show that
$\kappa$ itself does not deloop, is \emph{not} claimed here; see the end of
the sketch.)
\end{proposition}
\begin{proof}[Sketch]
As Lemma~\ref{l:diag}, $[\Delta\tilde\delta_F]=[F\circ K(2,2)]-2[F]$ in the
group $\hc{K_2}{K(B,4)}$ (pointwise addition through $e_4$; this group
structure is not in the library and must be built). Further lemmas: $K(g+h,m)$ is the pointwise sum (Lemma~\ref{l:dbl} with two
homomorphisms, so $\Phi_B$ is additive); $\cE(c;(G^*)_*)=\cE(c)\circ G$
(precomposition in the second variable), which with Lemmas~\ref{l:E2} and
j gives $\kappa\circ K(2,2)=\kappa;K(4,4)$; and \code{KFun}. Then
$[\Delta\tilde\delta_F]=\Phi_B(4g)-\Phi_B(2g)=\Phi_B(2g)=[\Delta(c_2;K(2g,4)_*)]$,
and Theorem~A. \emph{Converse (open).} An inhabitant of $\Unl f$ gives,
through \code{lemma9}, a cell $(b,h,k)$ at $x_0$ with $b$ and $h$ not
known to be $w_0$ and $f_1$; moving it to $(w_0,f_1)$ goes through the
singleton contractions of \code{Levels} (\code{s1}--\code{s3}) and the
coherence \code{coh}, and the pointing parts of the resulting paths
$f=D_p$ live in identity types of $\Om^2W$, not in a set (unlike
\code{Mult}, where $c=0$). Alternatively, $f=\lp\hat F$ is multiplicative by
\code{loop-mult}, but the pointing coherence at $q=\refl$ must then be
proved by $J$. Neither is needed for $\pi_4(S^3)$ at this stage; if
step~3 needs necessity, it will be stated and proved separately. (Once normalised, the conclusion
$\tilde\delta_F=0$ from its pointwise vanishing uses \code{pi-set}: $\Map(K_2,M_{K(B,4)})$ is
a set, and $\Om M_{K(B,4)}\simeq\Map(K_2,K(B,3))$ is connected because its
set of components $\hc{K_2}{K(B,3)}$ is trivial by (H3).)
\end{proof}

\section{Correspondence with the formal proof}\label{s:corr}
\begin{itemize}
\item \code{Z2} (new, built first): \code{Z2ab}, \code{xor} laws,
  \code{z2-set}, \code{rho}, \code{rho-hom}, \code{twoH}, \code{zeroH},
  \code{two-zero} (checked; \code{z2.opens}) (Definition~\ref{d:z2}, Lemma~\ref{l:two}).
\item \code{Doubling[A]}: \code{pi1} (Lemma~\ref{l:pi1}), \code{pmo-set}, \code{pmo-pi} (Lemma~\ref{l:pmo-pi}),
  \code{K0} (Lemma~\ref{l:K0}), \code{sq}, \code{ap-sq}, \code{dbl}
  (Lemma~\ref{l:dbl}); at $\Z/2$: \code{sq0}, \code{s}, \code{s'}
  (Corollary~\ref{c:sq0}).
\item \code{CrossTerm} (checked; \code{cross.opens}): \code{LMult.lmult}
  (\code{loop-mult} with its statement written), \code{S'.s'}
  (Corollary~\ref{c:sq0}, end); \code{k21}, \code{loop-k22}, \code{k21-sq},
  \code{ETwo.phi-two}, \code{ETwo.e-two} (Lemma~\ref{l:E2}), \code{ETwo.e-zero}
  ($\cE(0)$ pointwise $\refl$), \code{l2-mult} (Lemma~\ref{l:O2mult}),
  \code{kbar}, \code{Kill.kill}, \code{Kill.N}, \code{Kill.N-lev}
  (Theorem~\ref{t:kill}); in \code{Cross}: \code{f}, \code{Dp}, \code{Dhat},
  \code{tau-fun}, \code{tau-def}, \code{unev-ev}, \code{tau}, \code{tau-pt},
  \code{tau-formula} (Lemma~\ref{l:tau}; $e_2$, $e_4$ are the locked
  \code{Lev} ones), \code{Cj}, \code{delta}, \code{delta-def}
  (Definition~\ref{d:cross}); \code{cell-of} (Lemma~\ref{l:cell}),
  \code{f-sq}, \code{diag-null}, \code{diag-cls} (Lemma~\ref{l:diag}: an equality of
  \emph{classes} of diagonals), \code{cross}, \code{cross-cell} (Theorem~\ref{t:cross}: Theorem~A's
  injectivity upgrades it to an actual equality $\tilde\delta=0$ of bipointed
  maps, which \code{cell-of} turns into the paths $f=D_p$; no representative
  is chosen).
\item \code{DeloopR} (checked; \code{deloop.opens}): \code{cell0}, \code{rel},
  \code{unloop-contr} (Theorem~\ref{t:unloop}),
  \code{Fhat}, \code{gamma}, \code{BR} (Definition~\ref{d:BR}), \code{sq-f}, \code{omega-BR},
  \code{BR-trunc}, \code{BR-conn} (Theorem~\ref{t:BR}), \code{phi}, \code{transl}, \code{mu},
  \code{tor}, \code{pt-contr}, \code{tor-loops}, \code{tor-conn}, \code{tor-trunc} (Corollary~\ref{c:mu}).
\end{itemize}
Each declaration has an entry in \code{cross.opens} or \code{deloop.opens}.
\end{document}
